\documentclass[12pt]{article}
\usepackage{amsmath, amssymb, amsthm}
\usepackage[margin=1in]{geometry}

\newtheorem{theorem}{Theorem}
\newtheorem{lemma}[theorem]{Lemma}

\title{\textbf{Derivation D3: Aging Policies and Priority Key Dynamics}}
\author{Ashu Newatia \\ \textit{Architect}}
\date{\today}

\begin{document}

\maketitle

\section*{Part 1: Linear Aging and Static Keys}

Under the linear aging policy, a patient's priority dynamically increases as they wait. Let the emergency queue be formulated as a min-heap, where a mathematically lower value represents a higher clinical priority. 

Let $S_i$ be the initial severity score of patient $i$, $A_i$ be their arrival time, and $C > 0$ be the constant aging rate. The dynamic priority of patient $i$ at any current time $t \geq A_i$ is defined as:
\begin{equation}
    P_i(t) = S_i - C(t - A_i)
\end{equation}

\begin{theorem}
For any two waiting patients, their relative priority order under linear aging is strictly preserved over time $t$, allowing the queue to use a static, time-independent key without requiring structural updates.
\end{theorem}

\begin{proof}
Consider two patients, $i$ and $j$, residing in the queue at time $t$. Patient $i$ has strictly higher priority than patient $j$ if and only if $P_i(t) < P_j(t)$. Substituting the dynamic priority definition:
\begin{align*}
    S_i - C(t - A_i) &< S_j - C(t - A_j) \\
    S_i - Ct + C A_i &< S_j - Ct + C A_j
\end{align*}
Since the term $-Ct$ is applied equally to all elements in the queue at time $t$, it can be added to both sides of the inequality:
\begin{equation}
    S_i + C A_i < S_j + C A_j
\end{equation}
This inequality is completely independent of the variable $t$. Therefore, by assigning each patient a static key $K_i = S_i + C A_i$ strictly at the time of insertion, the min-heap will naturally preserve the exact relative ordering of the dynamic linear aging policy. Zero `decrease_key` updates are mathematically required while the patient remains in the queue.
\end{proof}

\section*{Part 2: Step-wise Escalation and Dynamic Update Bounds}

Under the step-wise escalation policy, a patient's severity score decreases by a discrete step value $\delta$ for every continuous interval $\Delta$ (e.g., 30 minutes) they remain untreated in the waiting room.

\begin{theorem}
Under step-wise escalation, the exact number of \texttt{decrease\_key} operations required for patient $i$ is strictly bounded by the discrete escalation steps achieved during their wait time.
\end{theorem}

\begin{proof}
Let patient $i$ arrive at time $A_i$ and depart the queue (either through treatment or leaving) at time $t_{end}$. Their total continuous wait time is $W_i = t_{end} - A_i$.

By the definition of the step-wise policy, a priority escalation occurs if and only if the current wait time $(t - A_i)$ is a strictly positive integer multiple of $\Delta$. Therefore, an escalation occurs at discrete time intervals:
\begin{equation}
    t_k = A_i + k\Delta \quad \text{for} \quad k \in \mathbb{Z}^+
\end{equation}
Because priority updates only alter the relative ordering of the heap at these exact discrete thresholds, structural modifications are only necessary at $t_k$. 

The total number of updates $U_i$ applied to patient $i$ over their total wait time $W_i$ is simply the number of complete $\Delta$ intervals that fit within $W_i$. This is expressed via the floor function:
\begin{equation}
    U_i = \left\lfloor \frac{W_i}{\Delta} \right\rfloor
\end{equation}
Because the queue only issues a \texttt{decrease\_key} operation exactly when a step boundary is crossed, the total number of updates invoked by the application perfectly matches the mathematical number of escalation steps $\lfloor W_i / \Delta \rfloor$. 
\end{proof}

\end{document}